\documentclass[10pt,a4paper]{amsart}
\usepackage[margin=19mm]{geometry}
\usepackage{amsmath,amssymb,mathtools}
\usepackage[scr=rsfs,cal=euler]{mathalfa}
\usepackage{xfrac}
\usepackage[unicode,hidelinks,bookmarks=false]{hyperref}
\newcommand{\ie}{i.e.\ }
\newcommand{\cf}{cf.\ }
\newcommand{\resp}{respectively\ }
\newcommand{\Subsection}{\subsection}
\providecommand{\nobreakdash}{\nobreak\hbox{-}}
\setlength{\emergencystretch}{2em}
\multlinegap0pt
\usepackage{amssymb}
\usepackage{xcolor}
\usepackage{bm}
\usepackage{tikz}
\newtheorem{lemma}{Lemma}
\newtheorem{theorem}{Theorem}
\title{Enhanced Asymptotic Analysis of Continuous-Time Markov Branching Systems: Revisiting Limiting Structural Theorems}
\author{Azam A Imomov, Sarvar B. Iskandarov, Jakhongir B. Azimov,, Hurshidjon Q. Jumaqulov}
\date{Source: arXiv:2604.01257v1 (CC BY 4.0)}
\begin{document}
\maketitle

\noindent\emph{Source and licence.} Enhanced Asymptotic Analysis of Continuous-Time Markov Branching Systems: Revisiting Limiting Structural Theorems. Version arXiv:2604.01257v1 (CC BY 4.0), distributed by arXiv under the Creative Commons Attribution 4.0 International licence, http://creativecommons.org/licenses/by/4.0/, as recorded in the arXiv licence metadata for this exact version and shown on the abstract page https://arxiv.org/abs/2604.01257v1. Full licence notice: \texttt{licenses/arxiv-2604.01257-CC-BY-4.0.txt}.\par\smallskip

\noindent\textit{Editorial scope.} Counted unit: the article's theorem IIBTh4 (source0.tex lines 845--862). The article's complete original proof of it is delivered with it (source0.tex lines 864--983, one \texttt{proof} environment of 120 lines). No mathematics is written, repaired or re-derived: the statement and the proof are the article's own lines, in the article's own notation, and no retained line is substituted or re-flowed.  Retained uncounted as the source-local context the proof needs: lines 226--228; lines 230--232; lines 468--470; lines 472--475; lines 490--492; lines 506--520; lines 516--518; lines 692--695; lines 772--774; lines 789--791; lines 797--807; lines 817--820.  Nothing was omitted from any retained span, so the package carries no omission record.  No retained line contains a citation, so the wrapper carries no bibliography and no bibliography entry.  No retained line contains a figure environment, an image-inclusion command or drawing code, so no image is delivered and the package declares no asset dependency.  Editorial formatting only, in addition: an AMS \texttt{amsart} preamble that restates the article's own declarations and macros used by the retained lines, namely the environments \texttt{lemma}, \texttt{theorem} on separate counters, together with the standard packages those lines need.  The retained environments print in this document as Lemma 1, Theorem 1 in order of appearance, whereas the article prints the same items with the numbering of its own full document.  The complete native source of the article as distributed in the arXiv e-print for this exact version is bundled unchanged as \texttt{sources/197646/source0.tex}.
\begin{equation}                            \label{IIB2.1}
    {\mathcal{P}_{i}(t;s)}=\left[F(t;s)\right]^{i}\exp\left\{\int_{0}^{t}h\left(F(u;s)\right)du\right\},
\end{equation}

\begin{equation}                           \label{IIB2.2}
    \frac{\partial{F(t;s)}}{\partial{t}}=f\left(F(t;s)\right),
\end{equation}

\begin{equation}                              \label{IIB3.1}
    \tau(t;s)=\frac{{(\nu t)}^{1/{\nu}}}{\mathcal{N}(t)}{{\left[1+\frac{\mathcal{M}(s)}{t}\right]}^{1/{\nu}}},
\end{equation}

\begin{equation}                              \label{IIB3.2}
    {{\mathcal{N}}^{\nu}}(t)\mathcal{L}\left(\frac{{(\nu{t})}^{1/{\nu}}}{\mathcal{N}(t)}\right)\to{1}
    \qquad {\text{as}} \quad {t\to\infty},
\end{equation}

\begin{equation}                              \label{IIB3.3}
    {{\lambda}_{\omega}}:=\lim_{x\to\infty}{x^{\nu}}\omega(x)<\infty.
\end{equation}

\begin{lemma}                  \label{IIBLem1}
    Let conditions $[{f_{\nu}}]$, $[{{\mathcal{L}}_{\nu}}]$ and $[{{\omega}_{\nu}}]$ be satisfied. Then
\begin{equation}                              \label{IIB3.4}
    \frac{1}{\Lambda(R(t;s))}-\frac{1}{\Lambda(1-s)}=\nu{t}+
    \int_{\tau(0;s)}^{\tau(t)}{\frac{\varepsilon(x)}{{x^{1-\nu}}\mathcal{L}(x)}dx}+
    \mathcal{O}\left(1/t\right)
\end{equation}
    as $t\to\infty$, where $\tau(t)={{(\nu{t})}^{1/\nu}}\big/{\mathcal{N}(t)}$ and $\mathcal{N}(t)$ is defined
    in {(\ref{IIB3.2})}. By incorporating the supplementary constraint {(\ref{IIB3.3})}, we directly derive the
    following fundamental identity:
\begin{equation}                              \label{IIB3.5}
    \frac{1}{\Lambda(R(t;s))}-\frac{1}{\Lambda(1-s)}=\nu{t}-\frac{1}{\nu}\ln\left[\Lambda(1-s)\nu{t}\right]+\rho(t;s),
\end{equation}
    where $\rho(t;s)=o(\ln{t})$ as $t\to\infty$ uniformly in $s\in[0,1)$.
\end{lemma}

\begin{equation}                              
    \frac{1}{\Lambda(R(t;s))}-\frac{1}{\Lambda(1-s)}=\nu{t}-\frac{1}{\nu}\ln\left[\Lambda(1-s)\nu{t}\right]+\rho(t;s),
\end{equation}

\begin{equation}                             \label{IIB3.19}
    q(t)=\frac{\mathcal{N}(t)}{{(\nu{t})}^{1/{\nu}}}\left(1+\frac{\ln{{a}_{0}}\nu{t}}{{{\nu}^{3}}t}
    \left(1+o(1)\right)\right) \qquad {\text{as}} \quad {t\to\infty},
\end{equation}

\begin{equation}                \label{IIB3.24}
    \mathcal{L}\left(\frac{1}{\phi(y)}\right)=\mathcal{L}\left(\frac{1}{y}\right)\left(1+K(y)\omega\left(\frac{1}{y}\right)\right).
\end{equation}

\begin{equation}                \label{IIB4.1}
    {{\lambda}_{\psi}}:=\lim_{x\to\infty }{x^{\delta}}\psi(x)<\infty.
\end{equation}

\begin{align}                \label{IIB4.2}
\left\{
\begin{array}{l}
    \displaystyle{{\lambda}_{\omega}}-\mathcal{L}(t)=\frac{{{\lambda}_{\omega}}}{\nu{t^{\nu}}}{L_{\omega}}(t)
    \qquad \hfill {\text{under assumption {(\ref{IIB3.3})},}}
\vspace{2.8mm}  \\
    \displaystyle{{\lambda}_{\psi}}-\ell(t)=\frac{{{\lambda}_{\psi}}}{\delta{t^{\delta}}}{L_{\psi}}(t)
    \qquad \hfill {\text{under assumption {(\ref{IIB4.1})},}}  \\
\end{array}
\right.
\end{align}

\begin{equation}                \label{IIB4.3}
    {C_{\textsf{L}}}-{\textsf{L}}(t)=\frac{{C_{\textsf{L}}}}{\delta{t^{\delta}}}+\mathcal{O}\left(\frac{1}{t^{\nu}}\right)
    \qquad {\text{as}} \quad {t\to\infty},
\end{equation}

\begin{theorem}                      \label{IIBTh4}
    Let Basic Axioms hold with $\gamma<0$ and let $\mu:=2\delta-\nu>0$. If ${C_{\textsf{L}}}=|\gamma|$
    in {(\ref{IIB4.3})}, then
\begin{equation}                \label{IIB4.4}
    {{e}^{\mathcal{T}(t)}}\mathcal{P}(t;s)=\mathcal{U}(s)\cdot\left(1+\zeta(t;s)\right),
\end{equation}
    where the limiting GF $\mathcal{U}(s)$ has the form
\begin{equation}                \label{IIB4.5}
    \mathcal{U}(s)=\exp\left\{\frac{1}{{(1-s)}^{|\gamma|}}+
    \int_{{1}/{(1-s)}}^{\infty}{\frac{|\gamma|-\textsf{L}(u)}{{u^{1-|\gamma|}}}du}\right\}
\end{equation}
    and
\[
    \zeta(t;s)=\frac{|\gamma|}{\delta\mu}\frac{{{\mathcal{N}}^{\mu}}(t)}{{(\nu{t})}^{{\mu}/{\nu}}}\left(1+\sigma(t;s)\right),
\]
    where $\mathcal{N}(\cdot)\in{{\mathcal{SV}}_{\infty}}$ is defined with the property {\eqref{IIB3.2}}
    and $\sigma(t;s)=\mathcal{O}\left({\ln{t}}/{t}\right)$ as $t\to\infty$ uniformly in $s\in[0,d],\;d<1$.
\end{theorem}

\begin{proof}
    Using the backward Kolmogorov equation {\eqref{IIB2.2}} in {\eqref{IIB2.1}}, we have
\begin{eqnarray}                              \label{IIB4.6}
    {{{e}^{\mathcal{T}(t)}}\mathcal{P}(t;s)}
& = & \exp\left\{{{\left[\tau(t)\right]}^{|\gamma|}}+\int_{0}^{t}{h\left(F(u;s)\right)du}\right\} \nonumber\\
\ \nonumber\\
& = & \exp \left\{{{\left[\tau(t;s)\right]}^{|\gamma|}}+\int_{s}^{F(t;s)}{\frac{h(u)}{f(u)}du}+\Delta(t;s)\right\},
\end{eqnarray}
    where $\Delta(t;s)={{[\tau(t)]}^{|\gamma|}}-{{[\tau(t;s)]}^{|\gamma|}}$. It is easy to see that
\[
    {{\left[\tau(t;s)\right]}^{|\gamma|}}=\frac{1}{{\left( 1-F(t;s) \right)}^{|\gamma|}}
    =\frac{1}{{(1-s)}^{|\gamma|}}+\int_{s}^{F(t;s)}{\frac{|\gamma|}{{(1-u)}^{1+|\gamma|}}du}.
\]
    Then we can rewrite {(\ref{IIB4.6})} as follows:
\begin{equation}                \label{IIB4.7}
    {{e}^{\mathcal{T}(t)}}\mathcal{P}(t;s)=\exp\left\{\frac{1}{{(1-s)}^{|\gamma|}}+\mathcal{B}(t;s)+\Delta (t;s)\right\},
\end{equation}
    where
\[
    \mathcal{B}(t;s)=\int_{s}^{F(t;s)}{\left[\frac{h(u)}{f(u)}+\frac{|\gamma|}{{(1-u)}^{1+|\gamma|}}\right]du}.
\]

    First we estimate $\Delta (t;s)$. According to our notation and
    using the asymptotic formula {\eqref{IIB3.5}}, we have
\begin{equation}                \label{IIB4.8}
    {{\left[\tau(t;s)\right]}^{|\gamma|}}={{\left[\nu{t}\mathcal{L}\left(\frac{1}{R(t;s)}\right)\right]}^{{|\gamma|}/{\nu}}}
    \left[1-\frac{|\gamma|{{\lambda}_{\omega}}}{{\nu}^{3}}\frac{\ln[\Lambda(1-s)\nu{t}]}{t}\left(1+o(1)\right)\right]
\end{equation}
    as $t\to\infty$. At the same time, from the previously established
    formula {\eqref{IIB3.1}}, we can immediately write that
\[
    R(t;s)=q(t)-q(t)\mathcal{M}(t;s),
\]
    where $q(t)=R(t;0)$ as before and $\nu{t}\mathcal{M}(t;s)\to\mathcal{M}(s)$ as $t\to\infty$. Therefore,
    since $q(t)\downarrow{0}$ as $t\to\infty$, by utilizing formula {\eqref{IIB3.24}} for $\phi(y)=R(t;s)$
    with $y=q(t)$ and $K(y)=\mathcal{M}(t;s)$, we can deduce that
\begin{equation}                \label{IIB4.9}
    \mathcal{L}\left(\frac{1}{R(t;s)}\right)=\mathcal{L}\left(\frac{1}{q(t)}\right)
    \left[1+\mathcal{M}(t;s)\omega\left(\frac{1}{q(t)}\right)\right].
\end{equation}
    Given our assumptions and according to asymptotic representation {\eqref{IIB3.19}} and
    first assertion of {\eqref{IIB4.2}}, it immediately becomes obvious that
\begin{align}                \label{IIB4.10}
\left\{
\begin{array}{l}
    \omega\left({1}\big/{q(t)}\right)=\mathcal{O}\left({q^{\nu}}(t)\right)=\mathcal{O}\left({1}\big/{t}\right)
    \qquad \hfill {\text{as}} \quad {t\to\infty,}
\vspace{2.8mm}  \\
    \mathcal{L}\left({1}\big/{q(t)}\right)={{\lambda}_{\omega}}+\mathcal{O}\left({1}\big/{t}\right)
    \qquad \hfill {\text{as}} \quad {t\to\infty.}\\
\end{array}
\right.
\end{align}
    Therefore, the utilization of relations {\eqref{IIB4.9}} and {\eqref{IIB4.10}} transforms equality {\eqref{IIB4.8}} into the form
\[
    {{\left[\tau(t;s)\right]}^{|\gamma|}}={{\left[\nu{t}\mathcal{L}\left(\frac{1}{q(t)}\right)\right]}^{{|\gamma|}/{\nu}}}
    \left[1-\frac{|\gamma|{{\lambda}_{\omega}}}{{\nu}^{3}}\frac{\ln[\Lambda(1-s)\nu{t}]}{t}+
    \mathcal{O}\left(\frac{\mathcal{M}(s)}{{t^{2}}}\right)\right].
\]
    Hence
\begin{eqnarray}              \label{IIB4.11}
    {\Delta(t;s)}
& = & {{\left[\nu{t}\mathcal{L}\left(\frac{1}{q(t)}\right)\right]}^{{|\gamma|}/{\nu}}}
    \frac{|\gamma|{{\lambda}_{\omega}}}{{\nu}^{3}}\frac{\ln\left[{\Lambda(1-s)}/{a_0}\right]}{t}
    \left(1+\mathcal{O}\left(\frac{1}{t}\right)\right) \nonumber\\
\ \nonumber\\
& = & \frac{|\gamma|C_{\mathcal{L}}^{1-{\delta}/{\nu}}{{\lambda}_{\omega}}}{{\nu}^2}
    \frac{\ln\left[{\Lambda(1-s)}/{a_0}\right]}{{(\nu{t})}^{{\delta}/{\nu}}}\left(1+\mathcal{O}\left(\frac{1}{t}\right)\right)
\end{eqnarray}
    as $t\to\infty$.

    Next, employing assumptions $[{f_{\nu}}]$ and $[{h_{\delta}}]$, we write
\begin{equation}                \label{IIB4.12}
    \mathcal{B}(t;s)=\int_{s}^{F(t;s)}{\frac{|\gamma|-\textsf{L}\left({1}\big/{(1-u)}\right)}{{{u}^{1+|\gamma|}}}du}
    =:\mathcal{B}(s)-\mathcal{J}(t;s),
\end{equation}
    where
\begin{equation}                \label{IIB4.13}
    \mathcal{B}(s)=\int_{{1}\big/{(1-s)}}^{\infty}{\frac{|\gamma|-\textsf{L}(x)}{{x^{1-|\gamma|}}}dx}
\end{equation}
    and
\begin{equation}                \label{IIB4.14}
    \mathcal{J}(t;s)=\int_{\tau(t;s)}^{\infty}{\frac{|\gamma|-\textsf{L}(x)}{x^{1-|\gamma|}}dx}.
\end{equation}
    At this stage, when deriving equations {\eqref{IIB4.13}} and {\eqref{IIB4.14}}, we utilized the substitution $1-u={1}/{x}$
    to simplify the subsequent transformation. According to {\eqref{IIB4.3}}, the integrand in equation {\eqref{IIB4.14}}
    asymptotically behaves as ${|\gamma|}/{\delta{x^{1+\mu}}}$ if $x\to\infty$. Consequently, it becomes necessary
    to analyze the integral $\int_{\tau(t;s)}^{\infty}{\left[|\gamma|/{\delta{x^{1+\mu}}}\right]dx}$ as $t\to\infty$.
    Then certainly that
\begin{equation}                \label{IIB4.15}
    \mathcal{J}(t;s)=\frac{|\gamma|}{\delta\mu}{{R}^{\mu}}(t;s)\left(1+o(1)\right)
    \qquad {\text{as}} \quad {t\to\infty}.
\end{equation}
    On the other hand, {\hyperref[IIBLem1]{Lemma}~\ref{IIBLem1}} implies
\begin{equation}                \label{IIB4.16}
    R(t;s)=\frac{\mathcal{N}(t;s)}{{(\nu{t})}^{{1}/{\nu}}}
    \left(1+\frac{\ln\left[\Lambda(1-s)\nu{t}\right]}{{{\nu}^{3}}t}\left(1+o(1)\right)\right)
\end{equation}
    as $t\to\infty$, where according to {\eqref{IIB4.9}},
    $\mathcal{N}(t;s)=\mathcal{N}(t)\left(1+\mathcal{O}\left({1}/{t}\right)\right)$. Then
    $\mathcal{J}(t;s)=\mathcal{O}\left( {1}/{{t^{\mu/\nu}}}\right)$. It is evident that both expressions $\Delta (t;s)$
    and $\mathcal{J}(t;s)$ approach zero. But we can confirm that the leading part of the tail term in the approximation
    of ${e^{\mathcal{T}(t)}}\mathcal{P}(t;s)$ by the limit function $\mathcal{U}(s)$ is estimated by $\mathcal{J}(t;s)$,
    since there is no doubt that $\mathcal{O}\left({1}/{t^{\mu/\nu}}\right)$ approaches zero much less rapidly
    than $\mathcal{O}\left({1}/{t^{\delta/\nu}}\right)$ does it as $t\to\infty$. It follows from {\eqref{IIB4.15}}
    and {\eqref{IIB4.16}}, that
\begin{equation}                \label{IIB4.17}
    \mathcal{J}(t;s)=\frac{|\gamma|}{\delta\mu}\frac{{{\mathcal{N}}^{\mu}}(t)}{{(\nu{t})}^{{\mu}/{\nu}}}
    \left(1+\frac{\mu}{{\nu}^{3}}\frac{\ln\left[\Lambda(1-s)\nu{t}\right]}{t}\left(1+o(1)\right)\right)
\end{equation}
    as $t\to\infty$. Therefore, by combining {\eqref{IIB4.7}}, {\eqref{IIB4.12}}, {\eqref{IIB4.17}},
    we conclude that
\[
    {e^{\mathcal{T}(t)}}\mathcal{P}(t;s)=\mathcal{U}(s)\exp\left\{\mathcal{J}(t;s)\right\}.
\]
    It suffices to observe that ${e^x}\sim{1+x}$ as $x\to{0}$, thereby establishing the convergence
    in {\eqref{IIB4.4}}, accompanied by the explicit remainder estimate given in the form of $\zeta(t;s)$.

    The theorem is proved completely.
\end{proof}

\end{document}
